% =============================================================================
% ARTIGO MODELO — QUALIS A1
% Classe principal: elsarticle (Elsevier)
% Classes alternativas: acmart (ACM) | sn-jnl (Springer Nature) | mdpi (MDPI)
% =============================================================================
% Para alternar:
%   ACM (conferência):     \documentclass[sigconf]{acmart}
%   ACM (journal):         \documentclass[acmsmall]{acmart} (ou acmlarge, acmtog)
%   Springer Nature:       \documentclass[sn-basic]{sn-jnl} (ou sn-nature, sn-mathphys-num, etc.)
%   MDPI:                  \documentclass[journal,article,submit]{Definitions/mdpi}
% =============================================================================
\documentclass[preprint,12pt,review]{elsarticle}

% --- Pacotes essenciais ---
\usepackage[brazilian,portuges]{babel}
\usepackage[utf8]{inputenc}
\usepackage[T1]{fontenc}
\usepackage{amsmath,amssymb}
\usepackage{graphicx}
\usepackage{hyperref}
\usepackage{booktabs}
\usepackage{float}
\usepackage[hyphens,brazilian]{url}
\usepackage{microtype}

% --- Metadados do PDF ---
\hypersetup{
  pdfauthor={Nome do Autor},
  pdftitle={Título do Artigo},
  pdfsubject={Área de Concentração},
  colorlinks=true,
  linkcolor=blue,
  citecolor=blue,
  urlcolor=blue
}

% --- Dados do artigo ---
\journal{Nome do Periódico (Qualis A1)}

\begin{document}

% ============================================================================
% CABEÇALHO
% ============================================================================
\begin{frontmatter}

\title{Título do Artigo: Subtítulo (se houver)}

\author[ufc]{Nome do Primeiro Autor\corref{cor1}}
\ead{email.primeiro@dominio.br}

\author[ufc]{Nome do Segundo Autor}
\ead{email.segundo@dominio.br}

\author[outro]{Nome do Terceiro Autor}
\ead{email.terceiro@dominio.br}

\cortext[cor1]{Autor correspondente.}

\address[ufc]{Universidade Federal do Ceará (UFC), Programa de Pós-Graduação
em Tecnologia Educacional, Fortaleza, CE, Brasil.}

\address[outro]{Instituição do Terceiro Autor, Departamento, Cidade, Estado, Brasil.}

\begin{abstract}
O resumo deve apresentar, de forma concisa e estruturada: contexto e problema
de pesquisa, objetivo, metodologia, principais resultados e conclusões. ABNT
NBR 6028:2021 recomenda 150 a 250 palavras para artigos científicos. Deve
permitir que o leitor compreenda o escopo e as contribuições do trabalho sem
necessitar do texto completo. Não incluir citações bibliográficas, siglas
isoladas ou abreviaturas não definidas. Para periódicos internacionais,
redigir em inglês (abstract); para nacionais, em português (resumo),
acompanhado de versão em inglês (abstract).
\end{abstract}

\begin{keyword}
Palavra-chave 1 \sep Palavra-chave 2 \sep Palavra-chave 3 \sep
Palavra-chave 4 \sep Palavra-chave 5
\end{keyword}

\end{frontmatter}

% ============================================================================
% 1. INTRODUÇÃO
% ============================================================================
\section{Introdução}

A introdução deve situar o leitor no tema da pesquisa, apresentando o contexto
mais amplo em que o problema se insere. Deve-se demonstrar a relevância do
tema por meio de dados empíricos, lacunas na literatura ou controvérsias
teóricas que justifiquem a realização do estudo \cite{autor2023}.

Recomenda-se que a introdução responda às seguintes questões:
\begin{itemize}
  \item O que já se sabe sobre o tema? (revisão de literatura breve)
  \item O que ainda não se sabe? (lacuna de pesquisa)
  \item Por que é importante preencher essa lacuna? (justificativa)
  \item Como o estudo foi conduzido? (breve descrição metodológica)
  \item Quais são as principais contribuições? (achados e implicações)
\end{itemize}

A última frase da introdução deve anunciar o objetivo do artigo e,
opcionalmente, a estrutura das seções seguintes.

% ============================================================================
% 2. REFERENCIAL TEÓRICO
% ============================================================================
\section{Referencial Teórico}

O referencial teórico fundamenta a pesquisa em teorias, conceitos e
resultados de estudos anteriores. Deve ser organizado por temas ou
categorias analíticas, não por autores individuais.

\subsection{Subtema 1}

Discussão do primeiro constructo teórico com citações diretas e indiretas.
Segundo \cite{autor2021}, a citação indireta (paráfrase) é preferível por
demonstrar domínio do conteúdo. Citações diretas curtas (até 3 linhas) são
incorporadas ao texto entre aspas duplas.

As citações diretas longas (4 ou mais linhas) devem ser destacadas em bloco
recuado a 4 cm da margem esquerda, em corpo 10, sem aspas:

\begin{citacao}
A citação direta longa deve ser apresentada em parágrafo independente,
recuado 4 cm da margem esquerda, sem aspas, em fonte tamanho 10 e
espaçamento simples entre linhas. A indicação do autor, ano e página
deve constar fora do bloco, no padrão autor-data (NBR 10520:2023).
\end{citacao}

\subsection{Subtema 2}

Discussão do segundo constructo teórico, estabelecendo conexões com o subtema
anterior. Nesta seção, é importante demonstrar como os conceitos dialogam
entre si e com o problema de pesquisa.

\subsection{Subtema 3}

Aprofundamento teórico complementar, incluindo a apresentação do modelo de
análise ou quadro conceitual que orienta a investigação.

% ============================================================================
% 3. METODOLOGIA
% ============================================================================
\section{Metodologia}

Esta seção descreve os procedimentos técnicos e científicos adotados para
alcançar os objetivos propostos. Deve ser suficientemente detalhada para
permitir a replicação do estudo.

\subsection{Classificação da Pesquisa}

A pesquisa classifica-se como:
\begin{itemize}
  \item \textbf{Quanto à natureza:} básica ou aplicada
  \item \textbf{Quanto aos objetivos:} exploratória, descritiva ou explicativa
  \item \textbf{Quanto à abordagem:} quantitativa, qualitativa ou mista
  \item \textbf{Quanto aos procedimentos:} experimental, bibliográfica,
    documental, estudo de caso, survey, pesquisa-ação, etc.
\end{itemize}

\subsection{Instrumentos e Procedimentos}

Descrição dos instrumentos de coleta de dados (questionários, entrevistas,
observação, análise documental) e dos procedimentos de aplicação.

\subsection{Análise dos Dados}

Métodos de análise empregados:

\begin{equation}
  \text{Índice} = \frac{\sum_{i=1}^{n} X_i}{n}
  \label{eq:indice}
\end{equation}

Descrição do tratamento estatístico ou da técnica de análise de conteúdo
utilizada. A Equação~\ref{eq:indice} exemplifica o cálculo utilizado.

\subsection{Aspectos Éticos}

Para pesquisas envolvendo seres humanos, informar o número do parecer de
aprovação do Comitê de Ética em Pesquisa (CEP). Para estudos com animais,
informar a aprovação da Comissão de Ética no Uso de Animais (CEUA).

% ============================================================================
% 4. RESULTADOS
% ============================================================================
\section{Resultados}

Apresentação objetiva dos achados da pesquisa, sem interpretação ou
discussão. Utilizar tabelas, figuras e gráficos para sintetizar dados.

\begin{table}[H]
\centering
\caption{Descrição da tabela com resultados principais.}
\label{tab:resultados}
\begin{tabular}{lccc}
\toprule
\textbf{Variável} & \textbf{Grupo A} & \textbf{Grupo B} & \textbf{p-valor} \\
\midrule
Indicador 1 & 45,2 $\pm$ 3,1 & 52,8 $\pm$ 2,9 & $< 0,05$ \\
Indicador 2 & 78,6 & 81,3 & $0,12$ \\
Indicador 3 & 12,4 $\pm$ 1,8 & 9,7 $\pm$ 1,5 & $< 0,01$ \\
\bottomrule
\end{tabular}
\end{table}

\begin{figure}[H]
\centering
% \includegraphics[width=0.8\textwidth]{figuras/resultados_grafico.png}
\caption{Gráfico de barras comparando os grupos antes e depois da intervenção.}
\label{fig:resultados}
\end{figure}

% ============================================================================
% 5. DISCUSSÃO
% ============================================================================
\section{Discussão}

Interpretação crítica dos resultados à luz do referencial teórico e de
estudos anteriores. Confrontar os achados com a literatura:
\begin{itemize}
  \item Resultados que confirmam estudos prévios
  \item Resultados que divergem da literatura, com hipóteses explicativas
  \item Limitações do estudo e vieses potenciais
  \item Implicações teóricas e práticas
\end{itemize}

% ============================================================================
% 6. CONCLUSÃO
% ============================================================================
\section{Conclusão}

Retomar o objetivo da pesquisa e sintetizar as principais contribuições.
Apresentar:
\begin{itemize}
  \item Síntese dos achados mais relevantes
  \item Contribuições teóricas e práticas
  \item Limitações do estudo
  \item Agenda para pesquisas futuras
\end{itemize}

A conclusão não deve introduzir dados ou argumentos novos. Deve refletir
diretamente sobre os objetivos anunciados na introdução.

% ============================================================================
% FINANCIAMENTO
% ============================================================================
\section*{Financiamento}
Este trabalho foi apoiado por [Agência de Fomento, Processo nº XXXX].

% ============================================================================
% AGRADECIMENTOS
% ============================================================================
\section*{Agradecimentos}
Os autores agradecem [instituição, laboratório, colaboradores, etc.].

% ============================================================================
% DECLARAÇÕES
% ============================================================================
\section*{Declarações}
\textbf{Conflito de interesses:} Os autores declaram não haver conflito de
interesses.

\textbf{Uso de IA generativa:} Os autores declaram que [descrever se e como
ferramentas de IA foram utilizadas na produção do artigo, conforme políticas
editoriais vigentes].

\textbf{Contribuição dos autores:}
\begin{itemize}
  \item Autor 1: conceituação, metodologia, investigação, escrita (original)
  \item Autor 2: análise formal, supervisão, escrita (revisão e edição)
  \item Autor 3: curadoria de dados, visualização
\end{itemize}

% ============================================================================
% REFERÊNCIAS
% ============================================================================
\section*{Referências}

% elsarticle: estilos disponíveis:
%   \bibliographystyle{elsarticle-num}        % Numérico
%   \bibliographystyle{elsarticle-harv}        % Harvard (autor-ano)
%   \bibliographystyle{elsarticle-num-names}  % Numerado com nomes
%
% ACM:  \bibliographystyle{ACM-Reference-Format}
% SN:   \bibliographystyle{sn-basic} (ou sn-nature, sn-vancouver, etc.)
% MDPI: \bibliographystyle{Definitions/mdpi}

\bibliographystyle{elsarticle-harv}
\bibliography{artigo_modelo_qualis_a1}

\end{document}
